\chapter{The Transverse-Field Ising Model}
\label{chap:tfim}

The transverse-field Ising model is the most important one-parameter reduction of
the transverse-field XY mother model.  It is simple enough to display the full
solution explicitly, but rich enough to exhibit the main structures of the
lecture notes: Jordan--Wigner fermionization, BdG quasiparticles, parity-sector
bookkeeping, Kramers--Wannier duality, order--disorder variables, and the
Majorana continuum limit.

In the present convention the Hamiltonian is
\begin{equation}
\label{eq:tfim-hamiltonian}
    H_{\rm TFIM}
    =
    -J\sum_{j=1}^{L}\sigma_j^x\sigma_{j+1}^x
    -h\sum_{j=1}^{L}\sigma_j^z .
\end{equation}
It is obtained from the XY chain by setting $\gamma=1$.  In most physical
applications one takes $J>0$ and $h\geq0$.  The full lattice solution, however,
also makes the second critical line at $h=-J$ visible.  We therefore keep signs
explicit when doing so is useful.

\section{Reduction from the XY mother model}
\label{sec:tfim-from-xy-mother}

Setting $\gamma=1$ in the XY-chain Hamiltonian gives
\begin{equation}
    H_{\rm XY}\big|_{\gamma=1}
    =
    -J\sum_j\sigma_j^x\sigma_{j+1}^x
    -h\sum_j\sigma_j^z,
\end{equation}
which is exactly \cref{eq:tfim-hamiltonian}.  The exact solution is therefore
obtained by specializing the XY-chain BdG data:
\begin{equation}
\label{eq:tfim-xi-delta-from-xy}
    \xi(k)=2(h-J\cos k),
    \qquad
    \Delta(k)=2\ii J\sin k
\end{equation}
in the standard BdG orientation used in the preceding chapter.  The
quasiparticle dispersion is
\begin{equation}
\label{eq:tfim-spectrum}
    E_{\rm TFIM}(k)
    =
    2\sqrt{(h-J\cos k)^2+J^2\sin^2 k}
    =
    2\sqrt{h^2+J^2-2hJ\cos k}.
\end{equation}
For $J>0$ and $h\geq0$, the minimum occurs at $k=0$ and the excitation gap is
\begin{equation}
\label{eq:tfim-gap-positive-field}
    \Delta_{\rm gap}=2\abs{h-J}.
\end{equation}
Thus the bulk gap closes at
\begin{equation}
\label{eq:tfim-critical-point}
    h=J.
\end{equation}
If negative fields are also retained, a second gap closing occurs at $h=-J$ and
$k=\pi$.

\begin{remark}[Why the TFIM is not merely a classical Ising chain]
The exchange term in \cref{eq:tfim-hamiltonian} is diagonal in the $\sigma^x$
basis, whereas the field term is diagonal in the $\sigma^z$ basis.  These two
parts do not commute.  The transition at $h=J$ is therefore a genuine quantum
phase transition controlled by the competition between ordering in the $x$
direction and quantum fluctuations generated by the transverse field.
\end{remark}

\section{Jordan--Wigner form and boundary sectors}
\label{sec:tfim-jw-form}

Using the Jordan--Wigner convention fixed earlier,
\begin{equation}
    c_j=
    \left(\prod_{\ell<j}\sigma_\ell^z\right)\sigma_j^+,
    \qquad
    c_j^\dagger=
    \left(\prod_{\ell<j}\sigma_\ell^z\right)\sigma_j^-,
\end{equation}
with $n_j=(1-\sigma_j^z)/2$, the open-chain or bulk fermionic Hamiltonian is
\begin{equation}
\label{eq:tfim-jw-bulk}
    H_{\rm TFIM}^{\rm bulk}
    =
    -J\sum_j
    \left(
        c_j^\dagger c_{j+1}+c_{j+1}^\dagger c_j
    \right)
    -J\sum_j
    \left(
        c_j^\dagger c_{j+1}^\dagger+c_{j+1}c_j
    \right)
    -h\sum_j(1-2n_j).
\end{equation}
This is a spinless $p$-wave superconducting chain with hopping and pairing of
equal magnitude.  Fermion number is not conserved, but fermion parity
\begin{equation}
\label{eq:tfim-fermion-parity}
    P_F=(-1)^N,
    \qquad
    N=\sum_j n_j,
\end{equation}
is conserved.

For a periodic spin chain, the Jordan--Wigner string turns the boundary term
into a parity-dependent fermionic boundary condition.  In a fixed parity sector
$P_F=p$, one imposes
\begin{equation}
\label{eq:tfim-parity-boundary-condition}
    c_{L+1}=-p\,c_1.
\end{equation}
Thus the even sector has anti-periodic fermions, while the odd sector has
periodic fermions:
\begin{equation}
\label{eq:tfim-sector-momenta}
    p=+1:
    \quad
    k=\frac{2\pi m+\pi}{L},
    \qquad
    p=-1:
    \quad
    k=\frac{2\pi m}{L}.
\end{equation}
The thermodynamic spectrum is insensitive to this distinction, but the exact
finite-size spectrum and the spin partition function require the correct parity
projection.

\section{BdG block and Bogoliubov diagonalization}
\label{sec:tfim-bdg-block}

For each non-self-conjugate momentum pair $(k,-k)$, define the Nambu spinor
\begin{equation}
    \Psi_k=
    \begin{pmatrix}
        c_k\\ c_{-k}^\dagger
    \end{pmatrix}.
\end{equation}
The BdG block is
\begin{equation}
\label{eq:tfim-bdg-block}
    \mathcal{H}_{\rm TFIM}(k)
    =
    \begin{pmatrix}
        2(h-J\cos k)&2\ii J\sin k\\
        -2\ii J\sin k&-2(h-J\cos k)
    \end{pmatrix}.
\end{equation}
Equivalently,
\begin{equation}
\label{eq:tfim-bdg-vector}
    \mathcal{H}_{\rm TFIM}(k)
    =
    \bm d_{\rm TFIM}(k)\cdot\bm\tau,
    \qquad
    \bm d_{\rm TFIM}(k)=
    \bigl(0,-2J\sin k,2(h-J\cos k)\bigr).
\end{equation}
The Bogoliubov angle may be chosen as
\begin{equation}
\label{eq:tfim-bogoliubov-angle}
    \cos\theta_k=\frac{h-J\cos k}{
    \sqrt{(h-J\cos k)^2+J^2\sin^2 k}},
    \qquad
    \sin\theta_k=\frac{J\abs{\sin k}}{
    \sqrt{(h-J\cos k)^2+J^2\sin^2 k}}.
\end{equation}
The diagonal many-body form is
\begin{equation}
\label{eq:tfim-diagonal-form}
    H_{\rm TFIM}^{(p)}
    =
    E_0^{(p)}+
    \sum_{k\in\mathcal{K}_{\vartheta_p}}
    E_{\rm TFIM}(k)\,\gamma_k^\dagger\gamma_k,
\end{equation}
where the allowed momentum set depends on the parity sector as in
\cref{eq:tfim-sector-momenta}.  Equivalently, after choosing one representative
from each pair $\{k,-k\}$,
\begin{equation}
\label{eq:tfim-pair-diagonal-form}
    H_k
    =
    E_{\rm TFIM}(k)
    \left(
        \gamma_k^\dagger\gamma_k
        +
        \gamma_{-k}^\dagger\gamma_{-k}
        -1
    \right).
\end{equation}
The BCS ground state in a fixed sector is the $\gamma=1$ specialization of the
XY-chain ground state.

\section{Phases and order parameters}
\label{sec:tfim-phases-order-parameters}

For $J>0$ and $h\geq0$, the two phases are:
\begin{equation}
\label{eq:tfim-phase-summary}
\begin{array}{ccl}
    h>J &:& \text{paramagnetic phase polarized along } z,\\[2pt]
    0\leq h<J &:& \text{ferromagnetic Ising phase ordered along } x.
\end{array}
\end{equation}
In the paramagnetic phase the field dominates and the ground state is smoothly
connected to
\begin{equation}
    \ket{\uparrow_z\uparrow_z\cdots\uparrow_z}.
\end{equation}
In the ferromagnetic phase, the thermodynamic ground space breaks the global
Ising symmetry
\begin{equation}
\label{eq:tfim-z2-symmetry}
    U_Z=\prod_{j=1}^{L}\sigma_j^z,
    \qquad
    U_Z\sigma_j^xU_Z^{-1}=-\sigma_j^x.
\end{equation}
The two symmetry-broken states are approximately the two $x$-polarized states at
$h=0$.

The exact spontaneous magnetization for $J>0$ and $0\leq h<J$ is
\begin{equation}
\label{eq:tfim-exact-order-parameter}
    m_x
    =
    \lim_{r\to\infty}
    \sqrt{\langle \sigma_j^x\sigma_{j+r}^x\rangle}
    =
    \left[1-\left(\frac{h}{J}\right)^2\right]^{1/8}.
\end{equation}
It vanishes for $h\geq J$.  The exponent $1/8$ is the two-dimensional Ising order-parameter exponent, reflecting the quantum--classical correspondence between the
one-dimensional TFIM and the two-dimensional classical Ising model.

\section{Kramers--Wannier duality}
\label{sec:tfim-kramers-wannier-duality}

The operator duality is defined on the dual lattice sites $j+1/2$ by
\begin{equation}
\label{eq:tfim-dual-operators}
    \mu_{j+\frac12}^z
    =
    \sigma_j^x\sigma_{j+1}^x,
    \qquad
    \mu_{j+\frac12}^x
    =
    \prod_{\ell\leq j}\sigma_\ell^z.
\end{equation}
These operators satisfy the same Pauli algebra on the dual lattice.  In
particular,
\begin{equation}
\label{eq:tfim-dual-field-term}
    \sigma_j^z
    =
    \mu_{j-\frac12}^x\mu_{j+\frac12}^x.
\end{equation}
Ignoring boundary subtleties for the moment, the Hamiltonian becomes
\begin{equation}
\label{eq:tfim-dual-hamiltonian}
    H_{\rm TFIM}(J,h)
    =
    -J\sum_j\mu_{j+\frac12}^z
    -h\sum_j\mu_{j-\frac12}^x\mu_{j+\frac12}^x.
\end{equation}
After a relabelling of the dual lattice and a spin-axis relabelling, this has the
same form as the original Hamiltonian with the couplings exchanged:
\begin{equation}
\label{eq:tfim-coupling-exchange}
    H_{\rm TFIM}(J,h)
    \longleftrightarrow
    H_{\rm TFIM}(h,J).
\end{equation}
The self-dual point is therefore
\begin{equation}
\label{eq:tfim-self-dual-point}
    h=J.
\end{equation}
Together with the exact solution, this confirms that the self-dual point is the
actual critical point for $J,h>0$.

The dual order parameter is a disorder operator in the original spin variables:
\begin{equation}
\label{eq:tfim-disorder-correlator}
    \langle \mu_{j+\frac12}^x\mu_{j+r+\frac12}^x\rangle
    =
    \left\langle
    \prod_{\ell=j+1}^{j+r}\sigma_\ell^z
    \right\rangle.
\end{equation}
Thus the ferromagnetic phase of the original spins maps to the disordered phase
of the dual spins, and the paramagnetic phase maps to the ordered phase of the
dual variables.

\section{Continuum Majorana limit}
\label{sec:tfim-continuum-majorana-limit}

For $J>0$ and $h\simeq J$, the gap closes near $k=0$.  Expanding
\cref{eq:tfim-bdg-block} at small $k$ gives
\begin{equation}
\label{eq:tfim-majorana-bdg-expansion}
    \mathcal{H}_{\rm TFIM}(k)
    \simeq
    2(h-J)\tau_z
    -2Jk\,\tau_y.
\end{equation}
The continuum mass and velocity are therefore
\begin{equation}
\label{eq:tfim-majorana-mass-velocity}
    m=2(h-J),
    \qquad
    v=2J.
\end{equation}
In real space, the low-energy Hamiltonian may be written in terms of two real
Majorana fields $\chi_R$ and $\chi_L$ as
\begin{equation}
\label{eq:tfim-majorana-field-theory}
    H_{\rm cont}
    =
    \frac{\ii v}{2}\int \dd x
    \left(
        \chi_R\partial_x\chi_R
        -
        \chi_L\partial_x\chi_L
    \right)
    +
    \ii m\int \dd x\,\chi_R\chi_L.
\end{equation}
The critical point $m=0$ is the Ising conformal field theory with central charge
\begin{equation}
\label{eq:tfim-central-charge}
    c=\frac12.
\end{equation}
The two signs of $m$ correspond to the two gapped phases.  In the fermionic
Kitaev-chain description, the sign change of the Majorana mass is the continuum
version of the topological transition at $\abs{\mu}=2\abs{t}$.

\section{Correlation functions from the Gaussian solution}
\label{sec:tfim-correlators-gaussian}

The fermionic ground state is Gaussian.  Therefore all fermionic correlators are
obtained from the two-point covariance matrix, and all spin correlators follow
from Wick's theorem after expressing Jordan--Wigner strings in terms of Majorana
operators.

For example, the transverse magnetization is local in fermions:
\begin{equation}
\label{eq:tfim-z-magnetization-fermions}
    \sigma_j^z=1-2c_j^\dagger c_j.
\end{equation}
By contrast, the Ising correlator contains a Jordan--Wigner string:
\begin{equation}
\label{eq:tfim-x-correlator-string}
    \sigma_i^x\sigma_j^x
    =
    \left(c_i^\dagger-c_i\right)
    \left[\prod_{\ell=i+1}^{j-1}(1-2n_\ell)\right]
    \left(c_j+c_j^\dagger\right),
    \qquad i<j.
\end{equation}
In Majorana language this becomes a Pfaffian of a finite antisymmetric matrix.
For translation-invariant infinite chains, the Pfaffian often reduces to a
Toeplitz determinant.  This is the computational bridge between the exact
Bogoliubov solution and the exact long-distance order parameter
\cref{eq:tfim-exact-order-parameter}.

\section{Summary}
\label{sec:tfim-summary}

The TFIM is the $\gamma=1$ reduction of the XY mother model.  Its exact solution
is controlled by the BdG dispersion
\begin{equation}
    E(k)=2\sqrt{h^2+J^2-2hJ\cos k}.
\end{equation}
For $J,h>0$, the gap closes at $h=J$, the same point selected by
Kramers--Wannier self-duality.  The phase $h>J$ is paramagnetic, while
$0\leq h<J$ has Ising ferromagnetic order in the $x$ direction.  The continuum
critical theory is a massless Majorana fermion, and the two gapped phases are
distinguished by the sign of the Majorana mass.
